
The rapid expansion of sequencing data in genomics and healthcare allows various opportunities and analytical challenges, particularly in studying bacterial exotoxins.
High-throughput sequencing has generated vast amounts of genomic data, including sequences encoding exotoxins that lack functional and target specificity annotations.
However, experimental research on exotoxins is labour-intensive and time-consuming, making it difficult to utilize this data.
As bacterial exotoxins play key roles in various therapeutics, such as pain management, bacterial infection treatment, and neurodegenerative disease research, there is a clear opportunity for implementing efficient computational approaches to accelerate research.

While machine learning has shown its potential for analyzing complex protein sequences, there have been no studies that aimed target-specific diversity of exotoxins so far. And when it comes to exotoxin type machine learning classification, the fourth type of bacterial exotoxins have never been included in the classification tasks.

To address this limitation, we propose ExoType and ExoTarget, novel predictive models designed to classify exotoxins by functional type and predict their target specificity (vertebrate vs. invertebrate).

ExoType and ExoTarget prediction models use a nonredundant data set of 1,067 sequences.
These models utilize protein embeddings as input generated by protein language models with T5 architecture, integrating sequence derived embeddings (MMseqs2) and structure based embeddings (AlphaFold2) to capture a more comprehensive representation of exotoxin features.

For the ExoType, a Support Vector Machine algorithm performs best in terms of differentiating exotoxins based on their functional type, that achieves Matthews Correlation Coefficients of 0.62. 
On the other hand, for ExoTarget, although all models (Random Forest, \gls{SVM}, \gls{KNN}, Logistic Regression) show robust performance in terms of \gls{MCC}, no single algorithm shows a statistically significant advantage over the others.


These results highlight the potential of protein language models and machine learning to improve research and discovery in this critical area, which has much potential to grow.